\documentclass[10pt]{article}
\providecommand{\FIGDIR}{fig}
\newcommand{\DIGESTDATE}{8 September 2026}
% ===================== PM Research Watch - shared preamble =====================
\usepackage[T1]{fontenc}
\usepackage[utf8]{inputenc}
\usepackage{amsmath,amssymb}
\usepackage{tgpagella}
\usepackage{tgheros}
\renewcommand{\sfdefault}{qhv}
\renewcommand{\ttdefault}{lmtt}
\usepackage{microtype}
\usepackage[english]{babel}

\usepackage[a4paper,top=17mm,bottom=17mm,left=17mm,right=17mm,
            headsep=6mm,footskip=10mm]{geometry}
\usepackage{graphicx}
\usepackage{xcolor}
\usepackage{colortbl}
\usepackage{tikz}
\usetikzlibrary{calc}
\usepackage[most]{tcolorbox}
\usepackage{titlesec}
\usepackage{fancyhdr}
\usepackage{booktabs}
\usepackage{array}
\usepackage{longtable}
\usepackage{enumitem}
\usepackage{multicol}
\usepackage{ragged2e}
\usepackage[hidelinks]{hyperref}

% ------------------------------------------------------------------ palette
\definecolor{Ink}    {HTML}{14262E}
\definecolor{Deep}   {HTML}{0F4C5C}
\definecolor{Teal}   {HTML}{1F7A8C}
\definecolor{Sky}    {HTML}{4FA3C4}
\definecolor{Sage}   {HTML}{6FA287}
\definecolor{Amber}  {HTML}{E3A008}
\definecolor{Clay}   {HTML}{B5643C}
\definecolor{Coral}  {HTML}{D1495B}
\definecolor{Violet} {HTML}{6A5B9E}
\definecolor{Slate}  {HTML}{4A6572}
\definecolor{Paper}  {HTML}{FBF9F5}
\definecolor{Mist}   {HTML}{EEEAE2}
\definecolor{Rule}   {HTML}{D9D3C9}

\pagecolor{Paper}
\color{Ink}

\hypersetup{colorlinks=true, linkcolor=Deep, urlcolor=Teal, citecolor=Teal,
            pdftitle={PM Research Watch}, pdfauthor={Automated daily digest}}

% ------------------------------------------------------------------ headings
\titleformat{\section}
  {\sffamily\bfseries\large\color{Deep}}{}{0pt}
  {}[{\vspace{1.2mm}\color{Amber}\titlerule[1.1pt]}]
\titlespacing*{\section}{0pt}{5.5mm}{2.8mm}

\titleformat{\subsection}
  {\sffamily\bfseries\normalsize\color{Teal}}{}{0pt}{}
\titlespacing*{\subsection}{0pt}{3.6mm}{1.4mm}

\setlist[itemize]{leftmargin=4.2mm, itemsep=1.1mm, topsep=1.3mm, parsep=0pt,
                  label=\textcolor{Teal}{\footnotesize$\blacktriangleright$}}

% ------------------------------------------------------------------ header/footer
\pagestyle{fancy}
\fancyhf{}
\renewcommand{\headrulewidth}{0pt}
\renewcommand{\footrulewidth}{0pt}
\fancyhead[L]{\sffamily\scriptsize\color{Slate}PM RESEARCH WATCH\;\textcolor{Rule}{\textbar}\;\DIGESTDATE}
\fancyhead[R]{\sffamily\scriptsize\color{Slate}Particulate matter monitoring \& health effects}
\fancyfoot[C]{\sffamily\scriptsize\color{Slate}\thepage}

% ------------------------------------------------------------------ components
% metric tile
\newcommand{\metric}[3]{%
  \begin{tcolorbox}[enhanced, colback=white, colframe=Rule, boxrule=0.5pt,
    arc=1.6mm, left=2.2mm, right=2.2mm, top=1.8mm, bottom=1.8mm,
    borderline west={1.7mm}{0pt}{#3}, nobeforeafter, valign=center,
    height=17mm, width=\linewidth]
    \RaggedRight
    {\sffamily\bfseries\fontsize{17}{17}\selectfont\color{#3}#1}\\[0.5mm]
    {\sffamily\fontsize{6.5}{7.4}\selectfont\color{Slate}#2}
  \end{tcolorbox}}

% pull-quote / key insight box
\newtcolorbox{keybox}[1][]{enhanced, colback=Mist, colframe=Mist,
  arc=1.6mm, left=3.2mm, right=3.2mm, top=2.4mm, bottom=2.4mm,
  borderline west={1.7mm}{0pt}{Amber}, #1}

% subtopic band
\newcommand{\band}[2]{%
  \vspace{3.2mm}\noindent
  \begin{tcolorbox}[enhanced, colback=#2, colframe=#2, arc=1.2mm,
    left=2.6mm, right=2.6mm, top=1.1mm, bottom=1.1mm, boxrule=0pt,
    width=\linewidth, nobeforeafter]
    {\sffamily\bfseries\small\color{white}#1}
  \end{tcolorbox}\vspace{1.6mm}}

% one paper entry: {short cite}{title}{journal}{doi}{body}
\newcommand{\paperentry}[5]{%
  \par\vspace{1.9mm}\noindent
  {\sffamily\bfseries\footnotesize\color{Ink}#1}\;%
  {\sffamily\fontsize{7}{8}\selectfont\color{Slate}#3}\\[0.4mm]
  {\footnotesize\itshape\color{Deep}#2}\\[0.6mm]
  {\footnotesize #4}\\[0.3mm]
  {\fontsize{6.6}{7.6}\selectfont\color{Slate}\texttt{doi:}\,\href{https://doi.org/#5}{\texttt{\detokenize{#5}}}}
  \par}

% caption style
\newcommand{\figcap}[1]{\par\vspace{1.0mm}
  {\fontsize{7.2}{8.6}\selectfont\color{Slate}\textsf{\textbf{Fig.}}\; #1\par}\vspace{1.5mm}}

\providecommand{\permil}{\textperthousand}

\setlength{\parindent}{0pt}
\setlength{\parskip}{0pt}
\renewcommand{\arraystretch}{1.18}


\begin{document}

% ==================================================================== MASTHEAD
\thispagestyle{empty}
\begin{tikzpicture}[remember picture, overlay]
  \fill[Deep] (current page.north west) rectangle
      ($(current page.north east)+(0,-67mm)$);
  \fill[Amber] ($(current page.north west)+(0,-67mm)$) rectangle
      ($(current page.north east)+(0,-68.4mm)$);
  \foreach \x/\y/\r in {22/12/0.9, 41/26/1.5, 63/9/0.7, 88/31/1.1, 112/17/1.9,
                        138/29/0.8, 159/13/1.3, 178/54/1.0, 196/21/1.6,
                        30/44/0.6, 74/57/1.2, 125/48/0.9, 168/6/0.8, 191/59/0.7,
                        52/18/1.0, 100/60/0.6, 148/20/1.1, 66/49/1.4, 205/45/0.8}
    \fill[Sky, opacity=0.30] ($(current page.north west)+(\x mm,-\y mm)$) circle (\r mm);
\end{tikzpicture}

\vspace*{4mm}
{\sffamily\fontsize{7.6}{9}\selectfont\color{Sky}%
 AUTOMATED DAILY LITERATURE DIGEST\;\textbullet\;PhD RESEARCH WATCH\par}
\vspace{2.2mm}
{\sffamily\bfseries\fontsize{27}{29}\selectfont\color{white}%
 Particulate Matter\par}
{\sffamily\fontsize{27}{29}\selectfont\color{Sky}%
 Monitoring \& Health Effects\par}
\vspace{5.0mm}
{\sffamily\fontsize{8.4}{10}\selectfont\color{white}%
 Issue 2026-09-08\;\;\textcolor{Amber}{\textbar}\;\;Entry window 8 September 2026\;\;%
 \textcolor{Amber}{\textbar}\;\;18 records screened in scope\par}
\vspace{18mm}

% -------------------------------------------------------------- metric strip
\noindent
\begin{minipage}[t]{0.190\linewidth}\metric{18}{RECORDS IN SCOPE\\(34 EXCLUDED)}{Deep}\end{minipage}\hfill
\begin{minipage}[t]{0.190\linewidth}\metric{7}{SHIPPED ON\\METADATA ONLY}{Teal}\end{minipage}\hfill
\begin{minipage}[t]{0.190\linewidth}\metric{4}{WITH A HUMAN\\HEALTH ENDPOINT}{Sky}\end{minipage}\hfill
\begin{minipage}[t]{0.190\linewidth}\metric{12}{ON THE SENSING /\\EXPOSURE AXIS}{Amber}\end{minipage}\hfill
\begin{minipage}[t]{0.190\linewidth}\metric{13}{NOT VIA THE\\PUBMED LEGS}{Clay}\end{minipage}

\vspace{6mm}

% -------------------------------------------------------------- signal of day
\section*{\sffamily\bfseries\color{Deep}Signal of the day}
\vspace{-1.5mm}{\color{Amber}\rule{\linewidth}{1.1pt}}\vspace{2.5mm}

\begin{keybox}
{\sffamily\bfseries\small\color{Deep}A third of the variance in indoor PM$_{2.5}$ sits
\emph{between} homes and survives adjustment for outdoor concentration, windows, cooking and
purifier use --- which is the quantitative case for per-dwelling measurement rather than
per-city inference.}\\[1.4mm]
\footnotesize
Forty-eight residences across eight Northern Chinese cities, paired indoor and outdoor
PM$_{2.5}$ at 5-minute resolution over seven consecutive days in each of four seasons,
aggregated to \textbf{31{,}094 valid hourly observations} of 32{,}256 scheduled
(\textit{Atmosphere}). The linear mixed-effects model carries a residence-level random
intercept and an AR(1) residual structure, and the two $R^2$ values are the finding:
\textbf{marginal $R^2=0.71$} for the fixed effects alone against \textbf{conditional
$R^2=0.91$} once the dwelling identity is admitted, with \textbf{ICC $=0.34$}. The fixed
effects behave as expected --- outdoor PM$_{2.5}$ $\beta=0.872$, window-open fraction
$0.548$, cooking $0.046$, purifier-on fraction $-0.902$, with continuous purifier operation
associated with \textbf{59.4\% lower} indoor PM$_{2.5}$ --- and the authors are careful to
call that last figure observational rather than causal. What matters for network design is
the residual third: after the outdoor field, the occupant behaviour and the building
covariates are all in the model, dwelling identity still carries a third of the variance,
and no ambient monitor or infiltration-factor lookup recovers it. Eight cities of Northern
China with coal-based district heating is a specific and unusually dirty regime, seven days
per season is short for a behavioural covariate, and window-open fraction and purifier use
are self-reported into the same model that they dominate.\par
\end{keybox}

% -------------------------------------------------------------- exec brief
\section*{\sffamily\bfseries\color{Deep}What changed today}
\vspace{-1.5mm}{\color{Amber}\rule{\linewidth}{1.1pt}}\vspace{2.5mm}

\footnotesize
\begin{itemize}
\item \textbf{Seven of eighteen records ship with no abstract in existence anywhere.} This is
new, and it is a property of same-day windowing rather than of the papers. Nine
Crossref-by-ISSN records carried no abstract in the journal-scoped payload; each was then
re-checked \emph{per DOI} against Crossref directly, against Europe PMC, against Semantic
Scholar and against the publisher page. All four legs returned nothing --- ScienceDirect and
Taylor \& Francis both refused the fetch outright. Seven are unambiguously in scope on the
title and ship as \texttt{Metadata only}, tier C, with the finding fields stating what is and
is not established; two more were rejected because without an abstract their scope cannot be
confirmed as particulate at all. \textbf{Crossref indexes a DOI faster than publishers deposit
the abstract, and a one-day entry window sits inside that gap.}

\item \textbf{Crossref-by-ISSN supplied 12 of 18 records; the PubMed spine supplied 5.} The
by-ISSN leg returned 59 raw and shipped 12 --- \textit{Aerosol Sci Technol}, \textit{J
Aerosol Sci} twice, \textit{Atmos Meas Tech}, \textit{Atmos Chem Phys}, \textit{Atmosphere}
three times, \textit{Atmos Pollut Res}, \textit{Indoor Environments}. Nine of those twelve
sit in journals the PubMed spine does not index at all. Two more (\textit{Environ Int},
\textit{Environ Pollut}) \emph{are} PubMed-indexed and still arrived here first, which is the
same deposit-lag asymmetry seen from the abstract side.

\item \textbf{The PubMed connector added nothing: the harvester's recall was complete.} Run
independently on both axes with an \texttt{[EDAT]} window, the connector returned
\textbf{19 health-axis} PMIDs --- the identical set the harvester fetched --- and
\textbf{2 sensing-axis}, a subset of the harvester's 3. First run in this archive where the
connector is a pure confirmation rather than a supplement.

\item \textbf{Consensus returned three hits and all three were already published here} ---
1, 2 and 5 August. Verified by resolving each title to its DOI and looking it up in
\texttt{seen.json}, not by inspecting publication dates, which closes the verification
weakness noted on 4 September. The dedup index did exactly what it exists to do.

\item \textbf{The trial-watch zero is topical, not a registry outage.} ClinicalTrials.gov
posted \textbf{1{,}135 updates registry-wide} on 8 September and \textbf{0} matching air
pollution or particulate matter. Yesterday's zero was a Labor Day artefact and could not
distinguish the two cases; today's no-term control separates them cleanly.
\texttt{trials.json} unchanged at 11.

\item \textbf{Every confidence interval in the issue comes from one paper.} The forest plot
holds four pooled estimates from a single systematic review, and its own heterogeneity
statistics are the reason to read it sceptically: PM$_{10}$ gives the largest risk ratio,
\textbf{1.27 (1.11--1.47)}, on \textbf{three studies at $I^2=99.5\%$}; PM$_1$ gives the only
heterogeneity-free estimate, \textbf{1.07 (1.05--1.09), $I^2=0.0\%$}, on \textbf{two}. The
review's title promises polypharmacy and it found \textbf{zero eligible studies} on it, and
none on household air pollution, older adults, or the Americas and Oceania.

\item \textbf{Two source legs were rate-limited rather than broken.} arXiv returned HTTP 429
where it has previously returned an unusable date filter, and OpenAlex returned 429 rather
than its usual 403 paywall message. Neither contributed. The OpenAlex leg remains broken for
a different reason and is still routed around.
\end{itemize}

% ==================================================================== ANALYTICS
\section*{\sffamily\bfseries\color{Deep}Corpus analytics}
\vspace{-1.5mm}{\color{Amber}\rule{\linewidth}{1.1pt}}\vspace{3mm}

\noindent\includegraphics[width=\linewidth]{\FIGDIR/f1_subtopics.png}
\figcap{Subtopic assignment of the 18 in-scope records. Only \textbf{six of the ten clusters
are occupied}, the narrowest spread in recent issues: \textit{Mechanistic toxicology},
\textit{Cardiovascular \& metabolic}, \textit{Neuro / mental health} and \textit{Respiratory
\& allergic} are all empty. \textit{Sensing, forecasting \& instrumentation} takes
\textbf{8 of 18} on its own. This is a source-leg effect --- Crossref-by-ISSN watches
instrumentation journals and fired hard today --- not a shift in the field.}

\vspace{2mm}
\noindent\includegraphics[width=\linewidth]{\FIGDIR/f2_design_endpoint.png}
\figcap{Left --- study architecture. The legend sums to 18 as \textbf{Metadata only 7,
measurement campaign 5, modelling / inventory 4, review / synthesis 1, observational (acute)
1}. The \textit{Metadata only} slice is the largest and is an artefact of the entry window,
not a description of the papers; every record in it is a real study whose abstract had not
deposited by the time the window closed. \textbf{Exactly one observational epidemiological
design appears in the whole issue.} Right --- health endpoint. \textbf{Fourteen of 18 carry
no health endpoint}, and all four that do are human, so the human-endpoint count is 4.}

\vspace{2mm}
\noindent\includegraphics[width=\linewidth]{\FIGDIR/f3_metric_geography.png}
\figcap{Left --- particle metric. \textit{PM (unspeciated) / emissions} is the largest bin at
7 of 18; only \textbf{one record reports a size-resolved or number metric} (Aitken-mode
particle number, and that one on metadata alone). \textit{Optical properties} at 3 is
unusually high and reflects the instrumentation skew. Right --- geographic provenance. No
real region dominates: Sub-Saharan Africa, Europe and China tie at 3 each. Three
Sub-Saharan records --- Accra, Ethiopia and a Western Sub-Saharan regional review --- ties
the archive's daily high, set on 6 August and matched on 3 September, but at \textbf{3 of 18
this is the largest \emph{share} the archive has recorded} (16.7\%, against 11.1\% and
7.9\%). The nominally largest bar, \textit{Global / multi-region} at 4, is a fallback: one
genuine multi-region meta-analysis, one bench study, and two metadata-only records with no
site stated.}

\vspace{2mm}
\noindent\includegraphics[width=\linewidth]{\FIGDIR/f5_forest.png}
\figcap{All four quantitative estimates in the corpus, on a common logarithmic axis.
\textbf{They are four pooled estimates from one systematic review, not four studies}, so the
panel shows a single meta-analysis's internal contrasts and nothing about between-study
agreement. Each is a random-effects adjusted risk ratio for multimorbidity per
5\,$\mu$g/m$^3$. Read the heterogeneity with them: $I^2$ is \textbf{93.0\%} for PM$_{2.5}$
(4 studies), \textbf{99.5\%} for PM$_{10}$ (3), \textbf{0.0\%} for PM$_1$ (2) and
\textbf{96.6\%} for NO$_2$ (4). The NO$_2$ interval touches 1.00. The visually widest
estimate, PM$_{10}$, is also the least interpretable.}

\vspace{2mm}
\noindent\includegraphics[width=\linewidth]{\FIGDIR/f4_heatmap.png}
\figcap{Subtopic $\times$ study-architecture cross-tabulation. The single densest cell is
\textit{Sensing} $\times$ \textit{Metadata only} at \textbf{4 records}, which is a statement
about publisher deposit latency in instrumentation journals rather than about study design.
Twelve cells are occupied out of sixty. The structural absence is the same one this watch
keeps recording: \textbf{no cell links an observational design to the sensing or exposure
clusters}. Nothing in this window measures particles with a novel or low-cost instrument
\emph{and} attaches a health outcome to what it measured; the Oklahoma cultivation-facility
study comes closest and stops at exposure characterisation.}

\vspace{2mm}
\noindent\includegraphics[width=\linewidth]{\FIGDIR/f6_lifecourse.png}
\figcap{Life-course windows addressed; counts are non-exclusive and only \textbf{5 of the 18
records} place anywhere on this axis. Three of the five carry a human health endpoint; the
other two --- the Ethiopian dormitories and the 48 Northern Chinese homes --- place by the
population whose exposure is characterised rather than by an outcome.}

\vspace{2mm}

\clearpage
% ==================================================================== DIGEST
\section*{\sffamily\bfseries\color{Deep}Paper-level digest}
\vspace{-1.5mm}{\color{Amber}\rule{\linewidth}{1.1pt}}\vspace{1mm}

{\fontsize{7.4}{9}\selectfont\color{Slate}Ordered by cluster. Each entry gives the
methodological core and the finding that matters, not an abstract paraphrase. Entries marked
\textsc{metadata only} had no abstract in existence on the entry date; for these the entry
states what the record establishes and what it does not, and asserts no finding.\par}

\band{Sensing, forecasting \& instrumentation\hfill 8 records}{Amber}

\paperentry{N China indoor Finf 48 homes 2026}
{Determinants of indoor PM$_{2.5}$ and infiltration factors across 48 dwellings}
{Atmosphere}
{Paired indoor/outdoor PM$_{2.5}$ at 5\,min in 48 residences across eight Northern Chinese
cities, seven consecutive days in each of four seasons, \textbf{31{,}094 valid hourly
observations} from 32{,}256 scheduled. A linear mixed-effects model with a residence random
intercept and AR(1) residuals gives \textbf{marginal $R^2=0.71$, conditional $R^2=0.91$,
ICC $=0.34$}: dwelling identity carries a third of the variance that the fixed effects ---
outdoor PM$_{2.5}$ ($\beta=0.872$), window-open fraction ($0.548$), cooking ($0.046$),
purifier-on fraction ($-0.902$) --- cannot reach. Continuous purifier operation associates
with \textbf{59.4\% lower} indoor PM$_{2.5}$, which the authors correctly decline to call
causal. Threat to inference: window and purifier use are self-reported and dominate the same
model, and coal-heated Northern China is a specific regime. For distributed sensing this is
the cleanest recent argument that per-dwelling measurement is not redundant with a good
ambient field.}
{10.3390/atmos17090874}

\paperentry{Europe BLH regression kriging 2026}
{Monthly boundary-layer height maps for Europe at 25\,km from radiosonde, satellite and ERA5}
{Atmos Meas Tech}
{Regression kriging over 2010--2020 combining radiosonde profiles with land-surface
temperature, topographic covariates and ERA5, at 25\,km and monthly means for 00 and 12 UTC.
Random forest beat linear regression and gradient boosting at both hours. Against ERA5 the
product is consistently sharper: winter spatial variability at 12 UTC gives
\textbf{RMSE $\leq$170\,m versus $\geq$250\,m}, summer temporal variability
\textbf{$\leq$260\,m versus 340\,m}. Performance \textbf{declines considerably at 00 UTC},
which is the honest limitation and the awkward one --- the nocturnal stable layer is exactly
where mixing-height error propagates hardest into surface concentration. Monthly means also
cannot serve episode-scale exposure work. Still, boundary-layer height is the dominant
meteorological term in almost every surface PM model, and a validated European product that
beats reanalysis is directly usable as a covariate.}
{10.5194/amt-19-5697-2026}

\paperentry{TROPOMI SO2 EAKF assimilation 2026}
{Assimilating TROPOMI SO$_2$ columns into a regional chemical transport model}
{Atmosphere}
{An Ensemble Adjustment Kalman Filter implemented in the MINNI regional CTM through DART,
assimilating Sentinel-5P TROPOMI SO$_2$-COBRA total columns over continental Europe for
August 2023, with a 20-member ensemble perturbing emissions and boundary conditions. The
filter constrained ensemble variance around power-plant and volcanic plumes; the monthly
domain-mean column correction was \textbf{2$\times$10$^{-5}$\,mol\,m$^{-2}$} with localised
maxima of \textbf{3.3$\times$10$^{-4}$}, and near-surface concentration corrections reached
\textbf{2.6\,$\mu$g\,m$^{-3}$}, validated against in-situ data. The stated limitations are
real ones: small ensemble, static vertical localisation, and the usual temporal fading of
initial-condition corrections. Relevance here is indirect but concrete --- SO$_2$ is a
sulfate precursor, and column-to-surface transfer skill is the same problem satellite-based
PM$_{2.5}$ products have to solve.}
{10.3390/atmos17090875}

\paperentry{SE India monsoon aerosol 2026}
{Seasonal aerosol optical properties and radiative effect at a rural Indian receptor site}
{Atmosphere}
{MICROTOPS-II sunphotometry from April 2021 to December 2023 at a rural site in Southeast
India, with trajectory-based source attribution (PSCF, CWT) and OPAC--SBDART radiative
transfer. Annual mean \textbf{AOD$_{500}$ 0.56 $\pm$ 0.22}, peaking pre-monsoon
(0.66 $\pm$ 0.19) and in winter (0.64 $\pm$ 0.23) and falling to 0.49 $\pm$ 0.21 in the
monsoon. \textbf{Angstrom exponent 1.30 $\pm$ 0.24} in the dry seasons against
\textbf{0.83 $\pm$ 0.37} in the monsoon marks the fine-to-coarse shift, corroborated by
negative spectral curvature. Precipitable water rose 2.00 to 4.35\,cm winter to monsoon,
which confounds the optical size retrieval through hygroscopic growth --- the paper says so.
Column optical depth is not surface mass, and a single rural receptor cannot be generalised;
its value is as a South Asian background record against which urban surface networks can be
referenced.}
{10.3390/atmos17090876}

\paperentry{PPDRTL leaf-image AQI 2026}
{Predicting roadside air quality from image texture of pollutant deposits on tree leaves}
{PLoS One}
{High-resolution images of roadside leaves at high, medium and low traffic-density sites in
Haryana, segmented with PSO, DPSO and FODPSO, reduced to contrast, entropy and standard
deviation features and regressed with ordinary linear regression onto Haryana State
Pollution Control Board reference AQI. The FODPSO variant reaches \textbf{$R^2$ up to 0.894}
with \textbf{RMSE 8.76--12.34\,$\mu$g\,m$^{-3}$} and reported site accuracies of 88.9--91.4\%.
The inference is weak for a reason worth naming precisely: the \textbf{pilot dataset is ten
days}, no held-out test set is described, and leaf deposition is a \emph{cumulative}
integrator being regressed onto an \emph{instantaneous} daily index, so the model is fitting
a mismatched target. Mixing an $R^2$ against a unitless AQI with an RMSE in
$\mu$g\,m$^{-3}$ compounds the confusion. Filed as a proxy-sensing curiosity, not a method.}
{10.1371/journal.pone.0326588}

\paperentry{SIBS inexpensive single-particle 2026}
{\textsc{metadata only} --- feasibility of a low-cost single-particle SIBS instrument}
{Aerosol Sci Technol}
{\textbf{No abstract existed in Crossref, Europe PMC, Semantic Scholar or on the publisher
page on the entry date}, and Taylor \& Francis refused the fetch. What the record
establishes: a feasibility study for an \emph{inexpensive} single-particle
spark-induced-breakdown-spectroscopy instrument has been published in the field's
instrumentation journal of record. What it does not establish: any detection limit, element
set, particle-size range, spark energy, hit rate or cost figure --- none of which can be
inferred from the title, and none of which is asserted here. This is the most directly
relevant record of the day to low-cost elemental speciation and is flagged for retrieval once
the abstract deposits.}
{10.1080/02786826.2026.2721653}

\paperentry{3D-LSS particle shape 2026}
{\textsc{metadata only} --- enhanced 3D light-scattering sensor for online particle shape}
{J Aerosol Sci}
{\textbf{No abstract retrievable on the entry date.} The title commits to two things: an
\emph{enhanced} three-dimensional light-scattering sensor for online aerosol shape
characterisation, and a validation against test aerosols. Neither the size range, the
angular sampling scheme, the sphericity metric nor the validation particle set is recoverable,
and none is asserted here. Shape is the property optical low-cost sensors silently assume
away when they convert scattering to mass, so an online shape channel is directly relevant;
the predecessor 3D-LSS work established the concept at 500\,nm--5\,$\mu$m, and what
``enhanced'' means against that baseline is precisely what the abstract would settle.}
{10.1016/j.jaerosci.2026.106860}

\paperentry{BC coating optics 2026}
{\textsc{metadata only} --- metal-oxide coatings and the optical properties of black carbon}
{J Aerosol Sci}
{\textbf{No abstract retrievable on the entry date.} The record establishes that the optical
effect of V$_2$O$_5$, NiO and Fe$_2$O$_3$ surface coatings on black carbon aggregates has
been characterised; it establishes no absorption enhancement factor, wavelength dependence,
coating thickness or morphology treatment, and none is asserted here. The reason to track it:
absorption-based BC measurement --- aethalometers, photoacoustic instruments and the
filter-based methods most networks run --- is systematically biased by the coating
enhancement factor, and the metal-oxide coatings specifically are the industrial and
brake-wear case rather than the usual organic or sulfate one.}
{10.1016/j.jaerosci.2026.106885}

\band{Exposure assessment \& modelling\hfill 4 records}{Teal}

\paperentry{Accra traffic intersections 2026}
{Meteorological drivers of PM$_{2.5}$ at nine Accra traffic hotspots}
{Environ Monit Assess}
{An Aeroqual Ranger portable monitor at nine traffic intersections in Accra across the
harmattan and non-harmattan seasons, against Ghana Meteorological Agency station data. Median
PM$_{2.5}$ was \textbf{54.90\,$\mu$g\,m$^{-3}$} at high-density sites, \textbf{34.26} at
commercial and \textbf{16.75} at low-density --- the cleanest category still exceeds the WHO
24\,h guideline of 15. Temperature correlated positively ($r=0.353$, $p<0.01$); rainfall,
wind speed and humidity negatively outside the harmattan, while in the harmattan wind speed
flipped weakly positive ($r=0.079$, $p=0.450$, i.e.\ null) consistent with dust resuspension.
The analysis is Pearson correlation only, with no multivariable adjustment and no traffic
counts, so ``the dominant factor remains localised vehicular emissions'' is a residual
argument rather than a measured one. A single portable instrument across nine sites also
gives no co-location error term.}
{10.1007/s10661-026-15889-8}

\paperentry{Xiamen chlorine NO2 uptake 2026}
{Chlorine-enhanced nocturnal NO$_2$ uptake and its nitrate yield in a coastal city}
{Atmos Chem Phys}
{Field observations in coastal Xiamen combined with a machine-learning attribution and a
multiphase box model, contrasting sea-land-breeze nights against non-SLB nights. HONO and
particulate nitrate both rose on SLB nights, with a mean NO$_2$ uptake rate constant of
\textbf{$k_{\mathrm{NO_2}}=9.70\times10^{-6}$\,s$^{-1}$}; chlorine ranked as the single most
important driver of the enhancement. Feeding that constant back into the box model resolved
the standing underestimation of both HONO and nitrate production, with nocturnal uptake
accounting for \textbf{83.9\% of HONO} and \textbf{46.4\% of nitrate formation}. The
machine-learning importance ranking is associational and cannot separate chlorine from the
marine air mass that carries it, and one coastal city in one season limits transfer.
Relevance is to secondary inorganic PM: nitrate is a large and poorly forecast fraction of
coastal urban PM$_{2.5}$.}
{10.5194/acp-26-12715-2026}

\paperentry{Brisbane Aitken-mode PNC 2026}
{\textsc{metadata only} --- long-term Aitken-mode particle number trends in Brisbane}
{Environ Int}
{\textbf{No abstract retrievable on the entry date} despite \textit{Environ Int} being a
PubMed-indexed journal --- the DOI indexed in Crossref before either the abstract or the
PubMed record deposited. The title commits to a long-term change in Aitken-mode particle
number concentration attributed across meteorology, traffic and biomass burning. Neither the
record length, the trend sign or magnitude, the instrument, nor the attribution method is
recoverable, and none is asserted here. Long records of sub-100\,nm number concentration are
rare and this one would bear directly on whether mass-based networks are tracking the
health-relevant fraction; flagged for retrieval.}
{10.1016/j.envint.2026.110514}

\paperentry{W Sub-Saharan Africa PM monitoring 2026}
{\textsc{metadata only} --- PM$_{2.5}$ and PM$_{10}$ monitoring capacity in Western
Sub-Saharan Africa}
{Atmos Pollut Res}
{\textbf{No abstract retrievable on the entry date.} The record establishes that a
challenges-and-perspectives treatment of PM$_{2.5}$ and PM$_{10}$ monitoring and estimation
across Western Sub-Saharan Africa has been published; whether it is a systematic review, a
narrative synthesis or a data-driven assessment cannot be determined from the metadata, and
no finding is asserted here. Western Sub-Saharan Africa has the sparsest reference network
of any populated region and is the canonical justification for low-cost sensor deployment, so
a current capacity assessment is worth retrieving --- particularly alongside the Accra
hotspot record in this same issue.}
{10.1016/j.apr.2026.103195}

\band{Occupational \& indoor\hfill 2 records}{Sage}

\paperentry{Oklahoma cannabis facility IAQ 2026}
{Seven-instrument occupational air-quality characterisation of a cannabis grow room}
{Ann Work Expo Health}
{Continuous monitoring from 10 July to 1 August 2023 in a 186\,m$^2$ HVAC-equipped medical
cannabis cultivation room holding roughly 900 plants on a 12:12 photoperiod. Seven
instruments across three spatial zones at 1.5\,m breathing height --- three Aranet4 NDIR
CO$_2$ sensors, two TSI AirAssure multi-parameter monitors carrying size-resolved PM, and two
PurpleAir PA-II-FLEX dual-channel counters --- produced \textbf{$\sim$113{,}000
quality-controlled records}. CO$_2$ tracked the light cycle, with means of
\textbf{979--1{,}086\,ppm} and maxima above \textbf{2{,}200}; nighttime elevation is
attributed to plant metabolism and enrichment rather than worker respiration, since workers
leave at lights-out. \textbf{Median PM$_{2.5}$ was low at 1.4--3.1\,$\mu$g\,m$^{-3}$} with
episodic activity-driven spikes. One facility, one room, one summer month, and no personal
sampling, so this is a determinants pilot rather than an exposure assessment --- but it is
the rare occupational study that runs consumer-grade and research-grade particle instruments
side by side.}
{10.1093/annweh/wxag074}

\paperentry{Ethiopia dormitory ventilation 2026}
{\textsc{metadata only} --- occupant-driven ventilation in naturally ventilated dormitories}
{Indoor Environments}
{\textbf{No abstract retrievable on the entry date.} The title commits to a field study of
occupant-driven ventilation and indoor exposure dynamics in naturally ventilated Ethiopian
university dormitories. The pollutants measured, the ventilation-rate method (tracer gas,
CO$_2$ decay or modelled), the number of rooms and the season are all unrecoverable and none
is asserted here. Occupant-controlled natural ventilation is the largest uncontrolled term in
indoor exposure models for warm low-resource settings, and measured air-change data from
Sub-Saharan Africa are scarce enough that the study is worth retrieving on those grounds
alone.}
{10.1016/j.indenv.2026.100192}

\band{Other clinical endpoints\hfill 1 record}{Violet}

\paperentry{Multimorbidity meta-analysis 2026}
{Pooled risk of multimorbidity per 5\,$\mu$g\,m$^{-3}$ of PM and NO$_2$}
{BMJ Public Health}
{Systematic review to 1 April 2026 across four databases (PROSPERO CRD42024613048); 22
studies included, 21 meta-analysed with random effects. Pooled adjusted risk ratios per
5\,$\mu$g\,m$^{-3}$: \textbf{PM$_{2.5}$ 1.15 (1.08--1.23)} from 4 studies at $I^2=93.0\%$;
\textbf{PM$_{10}$ 1.27 (1.11--1.47)} from 3 at $I^2=99.5\%$; \textbf{PM$_1$ 1.07
(1.05--1.09)} from 2 at $I^2=0.0\%$; \textbf{NO$_2$ 1.12 (1.00--1.26)} from 4 at
$I^2=96.6\%$. Nitrogen oxides, black carbon, SO$_2$, ozone and CO showed nothing. The
limitation is structural, not incidental: harmonising to a per-5\,$\mu$g\,m$^{-3}$ increment
discarded most of the 22 studies, leaving two to four per pollutant, and at
$I^2=93$--$99.5\%$ the PM$_{2.5}$ and PM$_{10}$ pooled values are averages over
non-exchangeable studies. \textbf{Zero eligible studies} on polypharmacy --- half the stated
aim --- or on household air pollution, older adults, the Americas or Oceania.}
{10.1136/bmjph-2025-004782}

\band{Reproductive \& developmental\hfill 1 record}{Coral}

\paperentry{France SUDI case-crossover 2026}
{Short-term air pollution and sudden unexpected death in infancy, national registry}
{Paediatr Perinat Epidemiol}
{\textbf{1{,}078 SUDI cases} from the French national registry, 2015--2022, with modelled
daily PM$_{10}$, PM$_{2.5}$, NO$_2$ and O$_3$ assigned at each infant's residence. A
time-stratified case-crossover design with conditional logistic regression, adjusted for
public holidays, influenza-like illness and daily mean temperature, examined lags to six
days. \textbf{PM exposure was consistently associated with higher odds of SUDI at lags 3--5
and across the last week of life}, strongest in cumulative and multi-pollutant models; NO$_2$
showed a trend only, and effects were larger in spring. The case-crossover design
self-matches on fixed confounders, which is its strength, but the abstract reports no point
estimates or intervals --- so the association's magnitude cannot be judged and this record
contributes nothing to the forest plot. Residence-level modelled exposure also ignores where
infants actually sleep and travel. SUDI remains an aetiologically open outcome; a short-term
environmental trigger is a plausible and previously little-tested component.}
{10.1111/ppe.70196}

\band{Burden, policy \& mitigation\hfill 2 records}{Clay}

\paperentry{Puyang PM2.5-O3 decade 2026}
{A decade of diverging PM$_{2.5}$ and ozone in a Chinese oilfield city}
{Preprints.org}
{\textbf{Preprint, not peer reviewed.} Regulatory monitoring records for Puyang, one of the
``2+26'' channel cities, 2015--2025. \textbf{PM$_{2.5}$ fell 79.6 to 39.3\,$\mu$g\,m$^{-3}$
while MDA8 ozone rose 123.0 to 189.6} --- the canonical Chinese divergence, here in a city
whose oil and gas VOC inventory makes it an unusually clean test of it. A positive
PM$_{2.5}$--O$_3$ correlation appeared only in the compound-pollution corner (O$_3 >
160$, PM$_{2.5} < 35$), with seasonal $r$ of 0.38 winter, 0.37 summer, 0.26 autumn, 0.21
spring --- all weak. Ozone production efficiency declined over the record, correlated with
PM$_{2.5}$ at $R=0.46$--$0.65$, and stayed above 7 monthly, which the authors read as
NOx-limited. OPE inferred from ambient monitoring rather than measured, no VOC speciation
reported, and a single city; the control recommendation rests on that inference.}
{10.20944/preprints202609.0349.v1}

\paperentry{Australian air health risk index 2026}
{\textsc{metadata only} --- community-specific weather and air health risk indices}
{Environ Pollut}
{\textbf{No abstract retrievable on the entry date}, despite \textit{Environ Pollut} being
PubMed-indexed --- the same Crossref-ahead-of-deposit lag seen elsewhere in this issue. The
title commits to community- and outcome-specific weather-and-air health risk indices built
with explainable machine learning across \textbf{66 Australian communities}. The pollutant
set, the health outcomes, the index construction and its skill are all unrecoverable and none
is asserted here. Outcome-specific and community-specific index construction is the
methodological direction that would replace one-size AQI banding, and 66 communities is
enough to test whether the community-specific part earns its complexity; flagged for
retrieval.}
{10.1016/j.envpol.2026.129117}

% ==================================================================== TABLE
\section*{\sffamily\bfseries\color{Deep}Corpus register}
\vspace{-1.5mm}{\color{Amber}\rule{\linewidth}{1.1pt}}\vspace{2.5mm}

{\fontsize{6.7}{8.4}\selectfont
\setlength{\tabcolsep}{2.4pt}
\begin{longtable}{@{}>{\RaggedRight}p{25mm}>{\RaggedRight}p{41mm}>{\RaggedRight}p{27mm}%
>{\RaggedRight}p{23mm}>{\RaggedRight}p{18mm}>{\RaggedRight}p{15mm}@{}}
\toprule
\textsf{\textbf{First author}} & \textsf{\textbf{Focus}} & \textsf{\textbf{Journal}} &
\textsf{\textbf{Design}} & \textsf{\textbf{Metric}} & \textsf{\textbf{Region}}\\
\midrule
\endfirsthead
\toprule
\textsf{\textbf{First author}} & \textsf{\textbf{Focus}} & \textsf{\textbf{Journal}} &
\textsf{\textbf{Design}} & \textsf{\textbf{Metric}} & \textsf{\textbf{Region}}\\
\midrule
\endhead
\multicolumn{6}{@{}l@{}}{\cellcolor{Mist}\textsf{\textbf{\textcolor{Deep}{Neuro / mental health}}}}\\[0.6mm]
\href{https://doi.org/10.1001/jamanetworkopen.2026.25015}{Dales} & Wildfire-sourced PM$_{2.5}$ and ED visits for migraine and other primary headache & \textit{JAMA Netw Open} & Case-crossover & PM$_{2.5}$ (wildfire) & Canada\\
\href{https://doi.org/10.1093/ageing/afag215}{Glover} & Umbrella meta-review of urban environmental exposures and cognitive health & \textit{Age Ageing} & Umbrella review / meta-analysis & PM$_{2.5}$ & Global\\
\href{https://doi.org/10.1016/j.nsa.2026.107019}{Goossens} & In utero air pollution exposure and pre-/postnatal brain development & \textit{Neurosci Appl} & Narrative review & PM$_{2.5}$ / PM$_{10}$ / PM$_{\mathrm{coarse}}$ & Global\\
\href{https://doi.org/10.1016/j.isci.2026.116617}{Ma} & PM$_{2.5}$-induced neurotoxicity and depression: exposure risks and mechanisms & \textit{iScience} & Narrative review & PM$_{2.5}$ + components & Global\\
\href{https://doi.org/10.3760/cma.j.cn112150-20250804-00752}{Pan} & Maternal PM exposure and offspring neurodevelopment: effects and mechanisms & \textit{Zhonghua Yu Fang Yi Xue Za Zhi} & Narrative review & PM (bulk) & China\\
\addlinespace[0.9mm]
\multicolumn{6}{@{}l@{}}{\cellcolor{Mist}\textsf{\textbf{\textcolor{Teal}{Cardiovascular \& metabolic}}}}\\[0.6mm]
\href{https://doi.org/10.3760/cma.j.cn112338-20251120-00843}{Huang} & PM$_{2.5}$ and dyslipidaemia in rural Chinese children; dietary effect modification & \textit{Zhonghua Liu Xing Bing Xue Za Zhi} & Cross-sectional (multistage) & PM$_{2.5}$ & China\\
\href{https://doi.org/10.32394/pe/218966}{Kazmierska} & Air pollution and cardiovascular disease: review of current data & \textit{Przegl Epidemiol} & Narrative review & PM$_{2.5}$ / PM$_{10}$ & Global\\
\href{https://doi.org/10.1007/s00125-026-06814-2}{Xi} & Air pollution, built environment and temperature vs. prediabetes phenotypes & \textit{Diabetologia} & Prospective cohort & PM$_{2.5}$ & Germany\\
\href{https://doi.org/10.1161/JAHA.125.046917}{Yang} & Short-term PM$_{2.5}$ and hospitalisation/ED use in atrial fibrillation & \textit{J Am Heart Assoc} & Case-crossover & PM$_{2.5}$ (wildfire-influenced) & USA\\
\addlinespace[0.9mm]
\multicolumn{6}{@{}l@{}}{\cellcolor{Mist}\textsf{\textbf{\textcolor{Sky}{Reproductive \& developmental}}}}\\[0.6mm]
\href{https://doi.org/10.1016/j.jes.2025.10.043}{Chen} & Prenatal air pollution, fetal growth trajectory phenotypes and child neurodevelopment & \textit{J Environ Sci (China)} & Prospective cohort & PM$_{2.5}$ / PM$_{10}$ & China\\
\href{https://doi.org/10.1016/j.ecoenv.2026.120566}{Huang} & Joint prenatal bisphenol + air-pollutant mixture exposure and preschool IQ & \textit{Ecotoxicol Environ Saf} & Prospective cohort (mixtures) & PM$_{2.5}$ / PM$_{10}$ & China\\
\href{https://doi.org/10.1007/s40572-026-00552-8}{Peterson} & Ambient air pollution, PAHs and pubertal development: critical review & \textit{Curr Environ Health Rep} & Critical review & PM$_{2.5}$ components & Global\\
\href{https://doi.org/10.1016/j.envres.2026.125313}{Xie} & Long-term PM$_{2.5}$ and timing of pubertal development in boys and girls & \textit{Environ Res} & Prospective cohort & PM$_{2.5}$ & China\\
\addlinespace[0.9mm]
\multicolumn{6}{@{}l@{}}{\cellcolor{Mist}\textsf{\textbf{\textcolor{Sage}{Respiratory \& allergic}}}}\\[0.6mm]
\href{https://doi.org/10.1007/s00484-026-03289-z}{Bayar Muluk} & Climate change and otolaryngology: airway disease impacts and emission reduction & \textit{Int J Biometeorol} & Narrative review & PM$_{2.5}$ & Global\\
\href{https://doi.org/10.4103/mgr.MEDGASRES-D-25-00164}{Liao} & Short-term air pollution, respiratory symptoms and lung function in asthma & \textit{Med Gas Res} & Panel study & PM$_{2.5}$ / PM$_{10}$ & China\\
\href{https://doi.org/10.3390/epidemiologia7040097}{Sanchez de Prada} & Ambient PM and asthma-coded admissions in children 0-3 y, Spain & \textit{Epidemiologia} & Ecological panel & PM$_{2.5}$ / PM$_{10}$ & Spain\\
\href{https://doi.org/10.1016/j.jhazmat.2026.143076}{Yang} & Allergenic vs. toxic risk of airborne seed hairs compared with ambient PM$_{2.5}$ & \textit{J Hazard Mater} & Field measurement & PM$_{2.5}$ + bioaerosol & China\\
\addlinespace[0.9mm]
\multicolumn{6}{@{}l@{}}{\cellcolor{Mist}\textsf{\textbf{\textcolor{Violet}{Mechanistic toxicology}}}}\\[0.6mm]
\href{https://doi.org/10.1016/j.jes.2025.12.074}{Ding} & PM$_{2.5}$ promotes vascular calcification via miR-27b-5p/HES1 signalling & \textit{J Environ Sci (China)} & In vitro + organoid + murine & PM$_{2.5}$ & China\\
\href{https://doi.org/10.1007/s00360-026-01698-5}{Maraschi} & Osmoregulatory disruption and metal bioaccumulation from settleable atmospheric PM & \textit{J Comp Physiol B} & Controlled exposure (ecotox) & Settleable APM & Brazil\\
\href{https://doi.org/10.1016/j.jes.2025.11.018}{Wu} & Macrophage NLRP3 inflammasome in real-ambient PM-induced hepatic glucose dysregulation & \textit{J Environ Sci (China)} & Real-ambient murine + co-culture & Real-ambient PM & China\\
\addlinespace[0.9mm]
\multicolumn{6}{@{}l@{}}{\cellcolor{Mist}\textsf{\textbf{\textcolor{Clay}{Occupational \& indoor}}}}\\[0.6mm]
\href{https://doi.org/10.3389/fmedt.2026.1796828}{Chen} & Nested suction instrument to cut intraoperative surgical-smoke PM$_{2.5}$ & \textit{Front Med Technol} & Bench / device experiment & PM$_{2.5}$ (surgical smoke) & China\\
\href{https://doi.org/10.1093/annweh/wxag055}{Cole} & Determinants of respirable crystalline silica exposure in Australian tunnelling & \textit{Ann Work Expo Health} & Observational exposure study & Respirable dust / RCS & Australia\\
\href{https://doi.org/10.1016/j.jhazmat.2026.143081}{Ding} & Stage-resolved hair Zn-isotope dynamics as a firefighting combustion-exposure tracer & \textit{J Hazard Mater} & Repeated-measures biomarker & Combustion-derived PM & China\\
\href{https://doi.org/10.3390/jox16040119}{Ramsden} & Xenobiotic hazards in aircraft cabin air & \textit{J Xenobiot} & Review + risk economics & Ultrafine particles & Global\\
\addlinespace[0.9mm]
\multicolumn{6}{@{}l@{}}{\cellcolor{Mist}\textsf{\textbf{\textcolor{Amber}{Sensing, forecasting \& instrumentation}}}}\\[0.6mm]
\href{https://doi.org/10.1016/j.jenvman.2026.130515}{Shi} & Adaptive spatiotemporal graph transformer for multi-site multi-step PM$_{2.5}$ forecasting & \textit{J Environ Manage} & Deep-learning model & PM$_{2.5}$ & China / India\\
\href{https://doi.org/10.3390/bios16070380}{Tsuchiya} & Seasonal and indoor-outdoor characteristics of fluorescent aerosol particles in offices & \textit{Biosensors} & Real-time sensor campaign & Bioaerosol / coarse & Japan\\
\addlinespace[0.9mm]
\multicolumn{6}{@{}l@{}}{\cellcolor{Mist}\textsf{\textbf{\textcolor{Amber}{Exposure assessment \& modelling}}}}\\[0.6mm]
\href{https://doi.org/10.1029/2025GH001621}{Solorzano-Araque} & Multiscale exposure modelling: LOTOS-EUROS + mobility surveys + particle toxicity & \textit{GeoHealth} & CTM + mobility integration & PM$_{2.5}$ / PM$_{10}$ & Colombia\\
\addlinespace[0.9mm]
\multicolumn{6}{@{}l@{}}{\cellcolor{Mist}\textsf{\textbf{\textcolor{Coral}{Burden, policy \& mitigation}}}}\\[0.6mm]
\href{https://doi.org/10.1002/hsr2.72886}{Karamzad} & Ischaemic heart disease burden in MENA 1990-2021; PM attributable fraction 28.2\% & \textit{Health Sci Rep} & GBD secondary analysis & Ambient PM & MENA\\
\href{https://doi.org/10.1371/journal.pone.0347989}{Rees} & HEPA filtration to reduce respiratory infection in care homes & \textit{PLoS One} & RCT process evaluation & Filtered PM & UK\\
\href{https://doi.org/10.1016/j.jes.2026.01.009}{Shan} & Policy-driven emission reductions and post-pandemic rebound, high-resolution inventories & \textit{J Environ Sci (China)} & Emission inventory & PM emissions & China\\
\addlinespace[0.9mm]
\multicolumn{6}{@{}l@{}}{\cellcolor{Mist}\textsf{\textbf{\textcolor{Slate}{Other clinical endpoints}}}}\\[0.6mm]
\href{https://doi.org/10.3390/dj14070440}{Lupita} & PM$_{2.5}$, preventive-care capacity and caries/edentulism burden across the EU & \textit{Dent J} & Ecological panel & PM$_{2.5}$ & EU-27\\
\href{https://doi.org/10.1016/j.jdin.2026.06.008}{Park} & Long-term PM exposure and risk of androgenetic alopecia & \textit{JAAD Int} & Nationwide cohort & PM$_{2.5}$ / PM$_{10}$ & South Korea\\
\href{https://doi.org/10.3171/2026.3.PEDS2615}{Sun} & Environmental exposure and area deprivation vs. paediatric hydrocephalus outcomes & \textit{J Neurosurg Pediatr} & Registry cohort + geospatial & PM$_{2.5}$ (satellite-derived) & USA\\
\addlinespace[0.9mm]
\bottomrule
\end{longtable}}


\vspace{4mm}
\begin{keybox}
{\sffamily\bfseries\footnotesize\color{Deep}Provenance and method}\\[1.2mm]
{\fontsize{7.2}{8.8}\selectfont
Entry window \textbf{8 September 2026}, a single day, run against \texttt{[Date - Entry]}.
The run began at 22:00 local, satisfying the standing rule that the current calendar day is
not opened before then; \texttt{last\_entry\_date} was 7 September, so exactly one day was
owed and no gap exists.
\textbf{PubMed} (NCBI E-utilities) returned \textbf{19 health-axis} and \textbf{3
sensing-axis} hits for 20 unique records. The \textbf{PubMed connector}, run independently on
both axes, returned \textbf{19} and \textbf{2} --- the identical health set and a subset of
the sensing set, so it added nothing; this is the first run in the archive where it is pure
confirmation. Seven of the 20 PubMed records had already shipped in earlier issues and were
dropped on PMID/DOI. \textbf{Europe PMC} (\texttt{CREATION\_DATE}) returned \textbf{1} raw,
a Preprints.org record neither PubMed leg carried, and it shipped. \textbf{Crossref-by-ISSN}
returned \textbf{59 raw} and shipped \textbf{12} --- two thirds of the issue, from
\textit{Aerosol Sci Technol}, \textit{J Aerosol Sci}, \textit{Atmos Meas Tech},
\textit{Atmos Chem Phys}, \textit{Atmosphere}, \textit{Atmos Pollut Res},
\textit{Indoor Environments}, \textit{Environ Int} and \textit{Environ Pollut}.
\textbf{arXiv} and \textbf{OpenAlex} both returned HTTP 429 and contributed nothing;
OpenAlex's created-date filter remains paywalled independently of the rate limit.
\textbf{Consensus} returned three hits, \textbf{all three already published here} on 1, 2 and
5 August --- verified by resolving each title to a DOI and querying \texttt{seen.json}, not
by inspecting publication dates. \textbf{Scholar Gateway} was queried once, for the SIBS
instrument paper; its corpus is Wiley-weighted and held no same-day Taylor \& Francis or
Elsevier deposit, so it is not a usable abstract source for records this fresh.
\textbf{ClinicalTrials.gov} posted \textbf{1{,}135 updates registry-wide} on 8 September and
\textbf{zero} matching air pollution or particulate matter --- a topical zero confirmed by
no-term control, unlike yesterday's federal-holiday zero. \texttt{trials.json} unchanged at 11.
\textbf{Thirty-four records were excluded with a logged reason}: ten duplicates already in
\texttt{seen.json}, and beyond those, ozone-only and greenhouse-gas endpoints, cigarette-smoke
toxicology, aquatic and materials chemistry, non-aerosol metrology, a correction notice, and
two records whose particulate scope could not be confirmed without an abstract.
\textbf{Seven in-scope records carry no abstract anywhere} --- checked per DOI against
Crossref, Europe PMC, Semantic Scholar and the publisher page --- and ship as
\texttt{Metadata only}, tier C. Their entries state what the metadata establishes and assert
no finding; treat them as pointers, and note that the corresponding rows of the register
carry no result.
The DOI gate passed \textbf{18 of 18 with zero warnings}. Three new design labels were added
to the plotting maps in the same edit as the records that use them; the geography map fired
no unmapped diagnostic. \textit{Global / multi-region} remains a fallback bucket that absorbs
unstated settings and bench work --- here 2 of its 4 are metadata-only records with no site.
Subtopic, design, particle-metric and geography labels are assigned from title, abstract and
keywords; assignment is single-label by dominant contribution and is a judgement, not a
controlled vocabulary. Effect estimates are transcribed verbatim; all four come from one
meta-analysis --- see the forest-plot caption.
All DOI links resolve to the publisher of record. Abstracts remain the copyright of their
respective publishers.\par}
\end{keybox}

\end{document}
